% !PRINTED 77
\weberpage{77}

%<<CONT
nicht in $\mathfrak{B}'$ enthalten, und also muss, wenn $\beta_0', \beta_1', \dots \beta_\mu'$ die
Indices von $b'$ sind, für jedes $x$ eine der Congruenzen

\weq{16}{
\begin{aligned}
x\beta_0' + \beta_0 &\equiv 0 \pmod{2} \\
x\beta_1' + \beta_1 &\equiv 0 \pmod{q_1} \\
&\cdots\cdots\cdots\cdots \\
x\beta_\mu' + \beta_\mu &\equiv 0 \pmod{q_\mu}
\end{aligned}
}

befriedigt sein.

Sind nun zwei Zahlen $\beta', \beta$ nicht beide durch eine Primzahl $q$ theilbar, so kann man immer über $x$ so verfügen, dass $x\beta' + \beta$ nicht durch $q$ theilbar wird; man braucht nur, wenn $\beta$ durch $q$ theilbar ist, $x$ durch $q$ nicht theilbar, und wenn $\beta$ durch $q$ nicht theilbar ist, $x$ durch $q$ theilbar anzunehmen. Man kann also $x$ auch so bestimmen, dass es gleichzeitig mehreren solchen Forderungen genügt. Es folgt also aus (16), dass entweder $\beta_0', \beta_0$ durch $2$ oder $\beta_1', \beta_1$ durch $q_1 \dots$ oder $\beta_\mu', \beta_\mu$ durch $q_\mu$ theilbar sein müssen; d.\,h. eine der Congruenzen (15) muss auch erfüllt sein, wenn $b$ irgend eine Zahl aus $\mathfrak{B}$ ist, gleichviel ob sie in $\mathfrak{B}'$ vorkommt oder nicht. Endlich folgt wieder aus (16), indem wir unter $b, b'$ zwei beliebige Zahlen aus $\mathfrak{B}$ verstehen, dass von den Congruenzen (15) eine und dieselbe für alle Zahlen aus $\mathfrak{B}$ bestehen muss. Dies können wir nun so zusammenfassen:

\gesperrt{2. Wenn die Periode $\eta$ nicht zu $\mathfrak{A}$ gehört, so muss es unter den Primzahlen (einschliesslich 2), die mehr als einmal in $m$ aufgehen, eine geben, $q$, so dass für jede Zahl $b$ aus $\mathfrak{B}$ der dem $q$ entsprechende Index $\beta$ durch $q$ theilbar ist.}

Wir betrachten jetzt die Gruppe $Q$, die nach der Bezeichnungsweise (5) des \S.\,18 durch die Congruenz

\weq{17}{Q \equiv 1 \pmod{\dfrac{m}{q}}}

definirt ist, und bezeichnen mit $P$ die dazu reciproke Gruppe.

Nach dem Satze 7., \S.\,19 ist jede Gruppe, die wie $\mathfrak{B}$ die Eigenschaft hat, dass der dem $q$ entsprechende Index $\beta$ durch $q$ theilbar ist, ein Theiler der Gruppe $P$, und folglich ist nach \S.\,12, 8. $Q$ ein Theiler von $\mathfrak{A}$, also $\mathfrak{A}$ nach \S.\,19, 8. ein nicht primärer Theiler von $\mathfrak{N}$, und damit sind wir zu dem Satze gelangt:

\gesperrt{3. Ist $\mathfrak{A}$ ein primärer Theiler von $\mathfrak{N}$, so gehört die Periode $\eta$ zu der Gruppe $\mathfrak{A}$.}
